\documentclass[10pt,conference]{IEEEtran}
\usepackage{kotex}
\usepackage{amsmath,amssymb,bm}
\usepackage{booktabs}
\usepackage{graphicx}
\usepackage{siunitx}
\usepackage{url}

\title{WIRIS 스테레오 영상과 Pixhawk 텔레메트리를 이용한\\
좌표계 일관 장거리 UAM 거리 측정 수식 모델}

\author{
\IEEEauthorblockN{저자명 TBD}
\IEEEauthorblockA{소속 TBD\\
Email: TBD}
}

\newcommand{\R}{\mathbb{R}}
\newcommand{\vect}[1]{\bm{#1}}
\newcommand{\norm}[1]{\left\lVert #1 \right\rVert_2}
\newcommand{\quat}{\otimes}
\newcommand{\frame}[1]{\mathcal{F}_{#1}}

\begin{document}
\maketitle

\begin{center}
\fbox{\parbox{0.92\linewidth}{\textbf{ARCHIVED DRAFT.} This document preserves an earlier stereo-center, reprojection-correction, and target-telemetry comparison proposal. It does not describe the current StereoDistance runtime. The current implementation emits Raw \texttt{left\_camera\_rdf} measurements only and has no measurement-correction or target ULog/telemetry comparison layer; use \texttt{docs/COORDINATE\_CONTRACT.md} as the authoritative contract.}}
\end{center}

\begin{abstract}
본 논문은 좌우 WIRIS 카메라 영상, YOLO 객체 검출, Pixhawk 자세, RTK/GNSS baseline 및 비행 후 ULog 기준 위치를 결합하여 장거리 UAM 또는 드론 표적의 거리와 3차원 위치를 추정하는 좌표계 일관 수식 모델을 제안한다.
제안 방법은 단순 disparity 기반 거리 산출을 최종 출력으로 사용하지 않고, 각 카메라의 bounding box 중심을 optical ray로 변환한 뒤 Pixhawk body FRD 및 world NED 좌표계 계약에 따라 회전시켜 두 ray 사이의 최단 선분 중점을 3차원 측정값으로 정의한다.
또한 bbox 중심 편향을 보정하기 위해 truth 재투영 residual을 screen-disparity cell별 robust median profile로 학습하고, lateral regression을 방지하기 위해 보수적인 range-only forward fusion 정책을 적용한다.
실험 결과는 추후 실제 UAM/드론 비행 ULog와 시뮬레이션 검증 결과를 삽입한다.
\end{abstract}

\begin{IEEEkeywords}
stereo vision, ray triangulation, UAM surveillance, Pixhawk, RTK, coordinate frame, YOLO, reprojection correction
\end{IEEEkeywords}

\section{서론}

장거리 UAM 또는 드론 감시에서 영상 기반 객체 검출은 표적 존재 여부를 빠르게 제공하지만, 단일 영상 bounding box만으로 실제 거리를 안정적으로 산출하기 어렵다.
특히 장거리에서는 작은 pixel 중심 오차가 큰 거리 오차로 증폭되며, 카메라 좌표계와 항법 좌표계를 혼동하면 축별 오차 해석이 물리적 의미를 잃는다.

본 연구의 목표는 현장 운용 가능한 스테레오 측정 도구인 StereoDistance를 기반으로, 영상 검출값과 Pixhawk/RTK 텔레메트리를 결합하는 수식 모델을 정리하는 것이다.
본 논문의 기여는 다음과 같다.

\begin{itemize}
    \item OpenCV optical frame, PX4 body FRD, world NED, stereo-center frame을 명시적으로 분리한 좌표계 계약을 제시한다.
    \item 좌우 YOLO bbox 중심을 optical ray로 변환하고 Pixhawk attitude를 반영하여 closest-ray triangulation으로 3차원 거리와 위치를 추정한다.
    \item ULog 또는 실시간 telemetry truth와 YOLO 측정 row를 시간 매칭하여 좌표 정렬 profile과 오차 metric을 계산하는 평가 절차를 정의한다.
    \item bbox 중심 residual을 재투영 기반 profile로 보정하되, lateral 축 회귀를 방지하는 보수적 forward/range 보정 정책을 제안한다.
\end{itemize}

\section{시스템 개요}

시스템은 좌우 WIRIS 카메라, 좌우 Pixhawk, GNSS/RTK baseline, YOLO 검출기, PyQt 기반 운용 GUI, 세션 기록 및 비행 후 분석 모듈로 구성된다.
좌우 WIRIS 카메라는 RTSP stream으로 영상을 제공하고, YOLO detector는 각 frame의 표적 bbox 중심을 산출한다.
Pixhawk는 MAVLink와 ROS2/uXRCE 경로를 통해 attitude, GPS/RTK, MAG 및 ULog 분석용 telemetry를 제공한다.

본 논문에서 최종 거리 원점은 좌우 카메라 optical center의 중점, 즉 stereo baseline center이다.
따라서 추정 range는 왼쪽 카메라 원점 기준이 아니라 stereo-center 기준으로 해석한다.

\section{좌표계 계약}

표~\ref{tab:frames}는 본 논문에서 사용하는 좌표계와 축 계약을 요약한다.
축 이름이 $x,y,z$로 저장된 경우에도, 해당 frame의 alias 계약이 확인되기 전까지는 물리 축으로 해석하지 않는다.

\begin{table}[t]
\centering
\caption{Coordinate frame and axis contracts}
\label{tab:frames}
\begin{tabular}{lll}
\toprule
Frame & Axis contract & Usage \\
\midrule
$\frame{c}$ & Right, Down, Forward & OpenCV optical ray \\
$\frame{b}$ & Forward, Right, Down & PX4 body FRD \\
$\frame{w}$ & North, East, Down & PX4 local/world NED \\
$\frame{s}$ & Right, Down, Forward & stereo-center optical \\
$\frame{t}$ & North, East, Down & ULog/RTK truth \\
\bottomrule
\end{tabular}
\end{table}

전방 장착 카메라의 고정 mount profile은 camera $+Z$ forward가 body $+X$ forward에 대응하고, camera $+X$ right가 body $+Y$ right에 대응하며, camera $+Y$ down이 body $+Z$ down에 대응한다고 둔다.
따라서 camera vector $\vect{v}_c=[x_c,y_c,z_c]^\top$와 body vector $\vect{v}_b=[x_b,y_b,z_b]^\top$의 변환은

\begin{equation}
\vect{v}_b = R_{cb}\vect{v}_c,
\quad
R_{cb} =
\begin{bmatrix}
0 & 0 & 1 \\
1 & 0 & 0 \\
0 & 1 & 0
\end{bmatrix}.
\label{eq:camera_to_body}
\end{equation}

역변환은

\begin{equation}
\vect{v}_c = R_{bc}\vect{v}_b,
\quad
R_{bc}=R_{cb}^{\top}.
\label{eq:body_to_camera}
\end{equation}

\section{Pixel-Ray 모델}

YOLO bbox 중심 pixel을 $\vect{p}=[u,v]^\top$라 하고, 카메라 내부 파라미터를

\begin{equation}
K =
\begin{bmatrix}
f_x & 0 & c_x \\
0 & f_y & c_y \\
0 & 0 & 1
\end{bmatrix}
\end{equation}

로 둔다.
왜곡 보정이 적용된 pixel 또는 live path에서 왜곡을 무시한 근사 pixel에 대해 normalized optical coordinate는

\begin{equation}
x = \frac{u-c_x}{f_x},
\quad
y = \frac{v-c_y}{f_y}
\end{equation}

이며, camera optical ray는

\begin{equation}
\vect{r}_c =
\frac{1}{\sqrt{x^2+y^2+1}}
\begin{bmatrix}
x \\ y \\ 1
\end{bmatrix}.
\label{eq:pixel_ray}
\end{equation}

여기서 $\vect{r}_c$는 OpenCV optical frame의 Right/Down/Forward 방향 벡터이다.

\section{Pixhawk 자세를 이용한 Ray 변환}

왼쪽 Pixhawk의 body-to-world quaternion을 $q_L$, 오른쪽 Pixhawk의 body-to-world quaternion을 $q_R$라 한다.
오른쪽 카메라 ray $\vect{r}_{R,c}$를 왼쪽 카메라 frame으로 표현하기 위해 먼저 오른쪽 camera ray를 오른쪽 body frame으로 변환하고, world NED frame으로 회전한 뒤, 왼쪽 body 및 camera frame으로 역변환한다.

\begin{equation}
\vect{r}_{R\rightarrow L,c}
=
R_{bc}
R(q_L)^{-1}
R(q_R)
R_{cb}
\vect{r}_{R,c}.
\label{eq:right_ray_to_left}
\end{equation}

여기서 $R(q)$는 quaternion $q$에 대응하는 회전행렬이다.
동일한 연산은 quaternion product로도 표현할 수 있다.
순수 quaternion $(0,\vect{v})$에 대해

\begin{equation}
\operatorname{rot}(q,\vect{v})
=
\left[
q \quat (0,\vect{v}) \quat q^{-1}
\right]_{xyz}
\end{equation}

로 둔다.

RTK 또는 GNSS baseline으로 얻은 오른쪽 카메라 원점의 world vector를 $\vect{B}_w$라 하면, 이를 왼쪽 camera frame translation $\vect{T}_{R/L,c}$로 변환한다.

\begin{equation}
\vect{T}_{R/L,c}
=
R_{bc}R(q_L)^{-1}\vect{B}_w.
\label{eq:baseline_camera}
\end{equation}

baseline norm을 외부에서 측정한 값 $B$로 고정해야 하는 경우,

\begin{equation}
\vect{T}_{R/L,c}
\leftarrow
B\frac{\vect{T}_{R/L,c}}{\norm{\vect{T}_{R/L,c}}}.
\end{equation}

\section{Closest-Ray Triangulation}

왼쪽 camera ray를 $\vect{d}_1$, 왼쪽 camera frame으로 변환된 오른쪽 camera ray를 $\vect{d}_2$, 오른쪽 camera 원점 translation을 $\vect{T}$라 한다.
두 ray는

\begin{equation}
L(s)=s\vect{d}_1,
\quad
R(t)=\vect{T}+t\vect{d}_2
\end{equation}

로 표현된다.
이상적인 경우 두 ray는 한 점에서 만나지만, 실제 bbox 중심 오차와 attitude 오차 때문에 skew line이 된다.
따라서 두 ray의 최단 선분 양 끝점을 구하고 그 중점을 측정점으로 정의한다.

다음 scalar를 정의한다.

\begin{equation}
a=\vect{d}_1^\top\vect{d}_1,
\quad
b=\vect{d}_1^\top\vect{d}_2,
\quad
c=\vect{d}_2^\top\vect{d}_2,
\end{equation}

\begin{equation}
\vect{w}_0=-\vect{T},
\quad
d=\vect{d}_1^\top\vect{w}_0,
\quad
e=\vect{d}_2^\top\vect{w}_0.
\end{equation}

분모

\begin{equation}
\Delta = ac-b^2
\end{equation}

가 0에 가까우면 두 ray가 거의 평행하므로 측정값을 버린다.
그렇지 않으면 closest parameters는

\begin{equation}
s^\star = \frac{be-cd}{\Delta},
\quad
t^\star = \frac{ae-bd}{\Delta}
\label{eq:closest_parameters}
\end{equation}

이다.
$s^\star \le 0$ 또는 $t^\star \le 0$이면 표적이 카메라 뒤쪽에 있는 해로 보고 폐기한다.
두 closest point는

\begin{equation}
\vect{p}_1=s^\star\vect{d}_1,
\quad
\vect{p}_2=\vect{T}+t^\star\vect{d}_2
\end{equation}

이며, 최종 왼쪽 camera frame 측정점은

\begin{equation}
\hat{\vect{p}}_{L,c}
=
\frac{\vect{p}_1+\vect{p}_2}{2}.
\label{eq:left_camera_point}
\end{equation}

ray 불일치 품질 지표는

\begin{equation}
g = \norm{\vect{p}_1-\vect{p}_2}
\label{eq:closest_gap}
\end{equation}

로 둔다.

\section{Stereo-Center 기준 거리}

왼쪽 camera frame의 측정점 $\hat{\vect{p}}_{L,c}$는 왼쪽 optical center 기준이다.
본 논문의 range 원점은 stereo-center이므로 오른쪽 camera translation의 절반을 빼서 stereo-center point를 얻는다.

\begin{equation}
\hat{\vect{p}}_{s,c}
=
\hat{\vect{p}}_{L,c}
-
\frac{1}{2}\vect{T}.
\label{eq:stereo_center_point}
\end{equation}

camera optical contract에 따라

\begin{equation}
\hat{\vect{p}}_{s,c}
=
\begin{bmatrix}
\hat{x}_{right} \\
\hat{y}_{down} \\
\hat{z}_{forward}
\end{bmatrix}.
\end{equation}

최종 stereo-center range는

\begin{equation}
\hat{r} = \norm{\hat{\vect{p}}_{s,c}}
=
\sqrt{
\hat{x}_{right}^2+
\hat{y}_{down}^2+
\hat{z}_{forward}^2
}.
\label{eq:range}
\end{equation}

\section{Rectified Stereo 보조 모델}

논문 설명과 진단 metric을 위해 rectified stereo 근사도 함께 둔다.
좌우 bbox 중심을 $(u_L,v_L)$, $(u_R,v_R)$라 하고 disparity를

\begin{equation}
\delta = |u_L-u_R|
\end{equation}

로 정의하면, baseline $B$와 focal length $f_x$에 대해 forward depth는

\begin{equation}
Z = \frac{f_xB}{\delta}.
\label{eq:rectified_depth}
\end{equation}

평균 pixel 중심

\begin{equation}
\bar{u}=\frac{u_L+u_R}{2},
\quad
\bar{v}=\frac{v_L+v_R}{2}
\end{equation}

을 사용하면

\begin{equation}
X_{right} = \frac{\bar{u}-\bar{c}_x}{f_x}Z,
\quad
Y_{down} = \frac{\bar{v}-\bar{c}_y}{f_y}Z.
\label{eq:rectified_lateral}
\end{equation}

단, 실제 live 측정의 기본값은 Eq.~\eqref{eq:closest_parameters}--\eqref{eq:range}의 3D ray triangulation이다.
Eq.~\eqref{eq:rectified_depth}는 disparity 품질 확인과 reprojection correction 설명에 사용되는 보조 모델이다.

\section{Truth 정렬 Profile}

YOLO 측정점 $\hat{\vect{p}}_{s,c}$와 ULog 또는 RTK truth $\vect{p}_t$는 서로 다른 frame에 있다.
따라서 축별 오차를 계산하기 전에 rigid alignment profile을 적용한다.
profile은 rotation $R_{st}$와 origin offset $\vect{o}_{t}$를 저장한다.

\begin{equation}
\hat{\vect{p}}_{t}
=
R_{st}\hat{\vect{p}}_{s,c}
+
\vect{o}_t.
\label{eq:alignment_apply}
\end{equation}

여기서 $\hat{\vect{p}}_{t}$는 PX4 local NED frame의 추정점이다.
truth를 camera/stereo-center frame으로 역변환할 때는

\begin{equation}
\vect{p}_{s,c}^{truth}
=
R_{st}^{\top}
\left(\vect{p}_{t}-\vect{o}_t\right)
\label{eq:alignment_inverse}
\end{equation}

를 사용한다.
profile 추정에는 최소 3개 이상의 비공선 matched sample이 필요하며, line-only trajectory는 회전 자유도가 퇴화하므로 최종 profile로 사용하지 않는다.

\section{시간 매칭}

YOLO 측정 row의 timestamp를 $\tau_i^{yolo}$, ULog 또는 telemetry truth timestamp를 $\tau_j^{truth}$라 한다.
UTC timestamp가 있는 경우 UTC 기준으로 먼저 매칭하고, 그렇지 않으면 boot/local timestamp와 offset $\Delta\tau$를 사용한다.

\begin{equation}
j^\star(i)
=
\arg\min_j
\left|
\tau_i^{yolo}
+\Delta\tau
-\tau_j^{truth}
\right|.
\label{eq:time_match}
\end{equation}

허용 최대 시간 차이를 $\tau_{max}$라 하면

\begin{equation}
\left|
\tau_i^{yolo}
+\Delta\tau
-\tau_{j^\star}^{truth}
\right|
\le
\tau_{max}
\end{equation}

를 만족하는 row만 matched sample로 사용한다.
ULog 기반 결과와 live telemetry fallback 결과는 같은 대표 metric으로 합치지 않고 truth source별로 분리 보고한다.

\section{BBOX 재투영 보정}

YOLO bbox 중심은 장거리에서 거리 오차를 크게 증폭시킨다.
본 논문은 truth point를 좌우 camera image plane으로 재투영하여 bbox 중심 residual을 학습한다.
truth point의 좌우 재투영 pixel을 $(u_L^{gt},v_L^{gt})$, $(u_R^{gt},v_R^{gt})$라 하고 YOLO bbox 중심을 $(u_L,v_L)$, $(u_R,v_R)$라 하면 residual은

\begin{equation}
\Delta\vect{b}
=
\begin{bmatrix}
\Delta u_L\\
\Delta v_L\\
\Delta u_R\\
\Delta v_R
\end{bmatrix}
=
\begin{bmatrix}
u_L^{gt}-u_L\\
v_L^{gt}-v_L\\
u_R^{gt}-u_R\\
v_R^{gt}-v_R
\end{bmatrix}.
\label{eq:bbox_residual}
\end{equation}

lookup feature는 평균 화면 위치와 raw disparity이다.

\begin{equation}
\bar{u}=\frac{u_L+u_R}{2},
\quad
\bar{v}=\frac{v_L+v_R}{2},
\quad
\delta=|u_L-u_R|.
\end{equation}

각 sample은 screen cell $(i,j)$와 disparity bin $k$로 배정된다.

\begin{equation}
i = \operatorname{bin}_u(\bar{u}),
\quad
j = \operatorname{bin}_v(\bar{v}),
\quad
k = \operatorname{bin}_d(\delta).
\end{equation}

cell별 residual profile은 robust median으로 정의한다.

\begin{equation}
\Delta\vect{b}_{ijk}
=
\operatorname{median}
\left\{
\Delta\vect{b}^{(n)}
\mid
n\in\mathcal{C}_{ijk}
\right\}.
\label{eq:cell_median}
\end{equation}

학습 sample 수가 부족하거나 lookup cell이 없거나 disparity가 유효하지 않으면 보정은 no-op으로 처리한다.

\subsection{Component Gain 보정}

좌우 horizontal residual을 평균 화면 이동 성분과 disparity 성분으로 분리한다.

\begin{equation}
\Delta u_{shift}
=
\frac{\Delta u_L+\Delta u_R}{2},
\quad
\Delta u_{disp}
=
\frac{\Delta u_L-\Delta u_R}{2}.
\end{equation}

vertical residual은 평균 vertical shift로 둔다.

\begin{equation}
\Delta v_{shift}
=
\frac{\Delta v_L+\Delta v_R}{2}.
\end{equation}

gain $g_u$, $g_d$, $g_v$를 적용한 effective residual은

\begin{equation}
\tilde{\Delta u}_L
=
g_u\Delta u_{shift}
+
g_d\Delta u_{disp},
\end{equation}

\begin{equation}
\tilde{\Delta u}_R
=
g_u\Delta u_{shift}
-
g_d\Delta u_{disp},
\end{equation}

\begin{equation}
\tilde{\Delta v}_L
=
\tilde{\Delta v}_R
=
g_v\Delta v_{shift}.
\end{equation}

현재 기본 보수 설정은

\begin{equation}
g_d=0.125,
\quad
g_u=0.0,
\quad
g_v=0.025.
\label{eq:default_gains}
\end{equation}

이다.
보정 bbox 중심은

\begin{equation}
\tilde{\vect{b}}
=
\vect{b}
+
\tilde{\Delta\vect{b}}
\end{equation}

로 계산하고, 보정된 중심을 다시 triangulation에 넣어 corrected point를 얻는다.

\subsection{Range-Only Forward Fusion}

완전한 재삼각측량 결과는 일부 궤적에서 lateral 축 회귀를 만들 수 있다.
따라서 기본 출력 정책은 raw lateral 축을 보존하고 corrected forward 축만 사용한다.

\begin{equation}
\hat{x}_{right}^{corr}
=
\hat{x}_{right}^{raw},
\quad
\hat{y}_{down}^{corr}
=
\hat{y}_{down}^{raw},
\quad
\hat{z}_{forward}^{corr}
=
\hat{z}_{forward}^{reproj}.
\label{eq:range_only_forward}
\end{equation}

이 정책은 lateral 축을 개선했다고 주장하지 않으며, bbox 중심 보정이 주로 range/forward error를 줄이는지를 검증 대상으로 둔다.

\section{평가 Metric}

정렬 profile이 적용된 PX4 local NED 추정점과 truth를 각각

\begin{equation}
\hat{\vect{p}}_{t}
=
\begin{bmatrix}
\hat{N}\\\hat{E}\\\hat{D}
\end{bmatrix},
\quad
\vect{p}_{t}
=
\begin{bmatrix}
N\\E\\D
\end{bmatrix}
\end{equation}

로 둔다.
축별 오차는

\begin{equation}
\vect{e}_{NED}
=
\hat{\vect{p}}_{t}
-
\vect{p}_{t}
\end{equation}

이며, position error는

\begin{equation}
e_{pos}
=
\norm{\vect{e}_{NED}}.
\label{eq:position_error}
\end{equation}

stereo-center 기준 truth range는

\begin{equation}
r^{truth}
=
\norm{
\vect{p}_{t}-\vect{o}_{t}
}
\end{equation}

이고, range error는

\begin{equation}
e_{range}
=
\hat{r}-r^{truth}.
\label{eq:range_error}
\end{equation}

camera-axis 비교가 필요한 경우 truth를 Eq.~\eqref{eq:alignment_inverse}로 stereo-center frame에 역변환한 뒤

\begin{equation}
\vect{e}_{FRD}
=
\begin{bmatrix}
\hat{z}_{forward}-z_{forward}^{truth}\\
\hat{x}_{right}-x_{right}^{truth}\\
\hat{y}_{down}-y_{down}^{truth}
\end{bmatrix}
\end{equation}

를 사용한다.
정렬 전 YOLO Right/Down/Forward와 truth North/East/Down을 직접 뺀 값은 mixed-frame 진단값일 뿐 최종 축별 정확도 지표로 사용하지 않는다.

RMSE는 matched sample set $\mathcal{M}$에 대해

\begin{equation}
\operatorname{RMSE}_{axis}
=
\sqrt{
\frac{1}{|\mathcal{M}|}
\sum_{i\in\mathcal{M}}
e_{axis,i}^{2}
}
\end{equation}

로 계산한다.
보정 profile 검증의 strict non-regression gate는 모든 validation trajectory에서

\begin{equation}
\operatorname{RMSE}_{axis}^{corr}
\le
\operatorname{RMSE}_{axis}^{raw},
\quad
axis\in\{Forward,Right,Down\}
\label{eq:non_regression}
\end{equation}

을 만족해야 한다.

\section{실험 설계}

실험 결과는 추후 삽입한다.
본 논문의 실험은 다음 두 범주로 분리한다.

\subsection{Simulation 검증}

Simulation 검증은 profile 학습용 trajectory와 validation trajectory를 분리한다.
학습 trajectory는 screen 좌/중/우, 상/중/하 및 여러 disparity/range bin을 채우도록 구성한다.
validation trajectory는 학습에 사용하지 않은 seed와 경로를 사용한다.
평가 표에는 raw, reprojection-corrected, range-only-forward 결과의 Forward/Right/Down RMSE, position RMSE, range RMSE를 기록한다.

\begin{table}[t]
\centering
\caption{Simulation validation result template}
\label{tab:simulation_results}
\begin{tabular}{lrrrr}
\toprule
Method & Fwd RMSE & Right RMSE & Down RMSE & Pos RMSE \\
\midrule
Raw & TBD & TBD & TBD & TBD \\
Reprojection full & TBD & TBD & TBD & TBD \\
Range-only forward & TBD & TBD & TBD & TBD \\
\bottomrule
\end{tabular}
\end{table}

\subsection{Hardware/Flight 검증}

실제 비행 검증은 WIRIS visible 8배줌 intrinsics, RTK/GNSS baseline, Pixhawk attitude, YOLO 측정 row, ULog truth를 함께 저장한 세션에서 수행한다.
ULog local NED source를 우선 사용하고, local source가 없을 때만 global/GPS-derived local NED fallback을 사용한다.
ULog 결과와 live telemetry fallback 결과는 분리 보고한다.

\begin{table}[t]
\centering
\caption{Flight validation result template}
\label{tab:flight_results}
\begin{tabular}{lrrrr}
\toprule
Truth source & Matched & Range MAE & Pos RMSE & Max error \\
\midrule
ULog local NED & TBD & TBD & TBD & TBD \\
ULog GPS-derived NED & TBD & TBD & TBD & TBD \\
Live telemetry fallback & TBD & TBD & TBD & TBD \\
\bottomrule
\end{tabular}
\end{table}

\section{논의}

본 수식 모델은 좌표계와 원점을 명확히 분리한다는 장점이 있다.
특히 camera optical Right/Down/Forward와 PX4 local North/East/Down을 직접 비교하지 않고, alignment profile 또는 명시적 회전 변환을 거친 뒤 오차를 계산한다.
또한 closest-ray gap $g$를 함께 기록하므로 bbox pair, attitude, baseline이 서로 일관되지 않는 상황을 거리값과 분리해 진단할 수 있다.

한계도 명확하다.
첫째, bbox 중심은 표적의 실제 3D 기준점과 일치하지 않을 수 있다.
둘째, 장거리에서는 작은 pixel residual이 큰 range residual로 증폭된다.
셋째, 본 모델의 기본 장착 가정은 좌우 WIRIS와 Pixhawk/GNSS가 동일한 rigid rig 축을 공유한다는 것이다.
boresight, lever arm, zoom별 optical center 이동이 큰 경우에는 별도의 extrinsic calibration 모델이 필요하다.

\section{결론}

본 논문은 WIRIS 스테레오 영상과 Pixhawk 텔레메트리를 결합하는 장거리 UAM 거리 측정 수식 모델을 제시하였다.
제안 모델은 pixel-to-ray 변환, Pixhawk attitude 기반 좌표 변환, closest-ray triangulation, stereo-center range 정의, ULog truth 정렬, bbox 재투영 보정을 하나의 좌표계 일관 pipeline으로 구성한다.
추후 연구에서는 실제 UAM/드론 비행 ULog와 정량 실험 결과를 추가하여 range/position 정확도와 보정 profile의 일반화 성능을 검증한다.

\section*{Acknowledgment}

TBD.

\begin{thebibliography}{1}

\bibitem{hartley2004multiple}
R. Hartley and A. Zisserman, \emph{Multiple View Geometry in Computer Vision}, 2nd ed. Cambridge University Press, 2004.

\bibitem{pixhawk}
PX4 Autopilot Documentation, ``PX4 User Guide,'' available online: \url{https://docs.px4.io/}.

\bibitem{opencv}
OpenCV Documentation, ``Camera Calibration and 3D Reconstruction,'' available online: \url{https://docs.opencv.org/}.

\bibitem{yolo}
J. Redmon et al., ``You Only Look Once: Unified, Real-Time Object Detection,'' in \emph{CVPR}, 2016.

\end{thebibliography}

\end{document}
